\chapter{Kansrekening en toevalsvariabelen}\label{ch:g11:prob}

De kansrekening modelleert experimenten met een onzekere uitkomst: een worp
met een dobbelsteen, een getrokken kaart, een gespeeld spel. Dit hoofdstuk
zet de woordenschat op eindige verzamelingen op, en voert daarna het centrale
personage van de hele kansrekening in: de \emph{\omterm{def:g11:prob:rv}{toevalsvariabele}} en haar
\emph{\omterm{def:g11:prob:expectation}{verwachtingswaarde}}. Die begrippen worden in \cref{ch:g12:randvar}
herhaald en verdiept, en de voorwaardelijke kansen volgen in
\cref{ch:g12:condprob}.

\section{Kans op een eindige verzameling}

\begin{definition}[Kansmodel]\label{def:g11:prob:model}
Een experiment met eindig veel uitkomsten wordt gemodelleerd door zijn
\emph{uitkomstenverzameling}\index{uitkomstenverzameling}
$\Omega = \{\omega_1, \dots, \omega_n\}$ (de verzameling van de uitkomsten)
samen met een \emph{\omterm{def:g10:proba:distribution}{kans}} $\P$ die aan elke uitkomst $\omega_i$ een getal
$p_i \geq 0$ toekent, met $p_1 + \dots + p_n = 1$. Een
\emph{gebeurtenis}\index{gebeurtenis} is een deelverzameling
$A \subseteq \Omega$, en
\[
\P(A) = \sum_{\omega_i \in A} p_i .
\]
Zijn alle uitkomsten even waarschijnlijk ($p_i = \frac1n$: het \emph{uniform
model}\index{uniform model}), dan geldt
\[
\P(A) = \frac{\text{aantal uitkomsten in } A}
{\text{aantal uitkomsten in } \Omega}.
\]
\end{definition}

\begin{example}\label{ex:g11:prob:twodice}
Gooi met twee onderscheidbare dobbelstenen: $\Omega$ is de verzameling van
de $36$ geordende paren $(i, j)$ met $1 \leq i, j \leq 6$, alle even
waarschijnlijk. De \omterm{def:g11:prob:model}{gebeurtenis} $S$ = ``de som is $7$'' bevat de $6$ paren
$(1,6)$, $(2,5)$, $(3,4)$, $(4,3)$, $(5,2)$, $(6,1)$, dus is
$\P(S) = \frac{6}{36} = \frac16$. De \omterm{def:g11:prob:model}{gebeurtenis} ``de som is $12$'' bevat
alleen $(6,6)$: \omterm{def:g10:proba:distribution}{kans} $\frac{1}{36}$. De \emph{juiste} \omterm{def:g11:prob:model}{uitkomstenverzameling}
kiezen --- geordende paren, geen ongeordende --- is precies wat het \omterm{def:g11:prob:model}{uniform
model} bruikbaar maakt.
\end{example}

\begin{proposition}[Rekenregels voor kansen]\label{prop:g11:prob:rules}
Voor gebeurtenissen $A, B \subseteq \Omega$ geldt
\[
\P(\emptyset) = 0, \qquad \P(\Omega) = 1, \qquad
\P(\bar A) = 1 - \P(A),
\]
\[
\P(A \cup B) = \P(A) + \P(B) - \P(A \cap B),
\]
waarbij $\bar A$ het \emph{\omterm{def:g10:proba:operations}{complement}} van $A$ is (alle uitkomsten die niet
in $A$ zitten). In het bijzonder is $\P(A \cup B) = \P(A) + \P(B)$ wanneer
$A$ en $B$ \emph{\omterm{def:g10:proba:operations}{onverenigbaar}} zijn ($A \cap B = \emptyset$).
\end{proposition}

\begin{proof}
De lege \omterm{def:g11:prob:model}{gebeurtenis} bevat geen uitkomst (lege som), en $\Omega$ bevat ze
allemaal (totale som $1$). Elke uitkomst zit in precies één van $A$ en
$\bar A$, dus is $\P(A) + \P(\bar A) = 1$. Voor de \omterm{def:g10:numbers:interunion}{vereniging}: de $p_i$
optellen over $A$ en daarna over $B$ telt de uitkomsten van $A \cap B$ twee
keer; $\P(A \cap B)$ aftrekken herstelt die dubbele telling.
\end{proof}

\begin{example}
In een klas volgt $60\%$ muziek ($M$), doet $45\%$ aan sport ($S$), en doet
$25\%$ allebei. Het aandeel dat minstens één van beide doet, is
$\P(M \cup S) = 0.60 + 0.45 - 0.25 = 0.80$, zodat $20\%$ geen van beide doet
(\omterm{def:g10:proba:operations}{complement}).
\end{example}

\section{Toevalsvariabelen}

\begin{definition}[Toevalsvariabele, kansverdeling]\label{def:g11:prob:rv}
Een \emph{toevalsvariabele}\index{toevalsvariabele} op $\Omega$ is een
\omterm{def:g11:func:function}{functie} $X$ die aan elke uitkomst een getal toekent. Haar
\emph{kansverdeling}\index{kansverdeling} is de tabel van haar mogelijke
waarden $x_1, \dots, x_k$ samen met hun kansen
\[
\P(X = x_i) = \P\bigl(\{\omega : X(\omega) = x_i\}\bigr),
\qquad
\sum_{i=1}^k \P(X = x_i) = 1 .
\]
\end{definition}

\begin{example}\label{ex:g11:prob:game}
Een spel: betaal $2$ euro, gooi met één eerlijke dobbelsteen, en ontvang de
getoonde waarde wanneer die minstens $5$ is, en anders niets. Zij $G$ de
\emph{winst} (ontvangen min betaald). De uitkomsten $1, 2, 3, 4$ geven
$G = -2$; de uitkomst $5$ geeft $G = 3$; de uitkomst $6$ geeft $G = 4$. De
\omterm{def:g11:prob:rv}{kansverdeling} van $G$:
\begin{center}
\begin{tabular}{c|ccc}
$g$ & $-2$ & $3$ & $4$\\
\hline\rule{0pt}{11pt}%
$\P(G = g)$ & $\dfrac46$ & $\dfrac16$ & $\dfrac16$
\end{tabular}
\end{center}
\end{example}

\begin{method}[Een kansverdelingstabel opstellen]\label{met:g11:prob:table}
Noem de mogelijke waarden van $X$; verzamel voor elke waarde de uitkomsten
die haar voortbrengen en tel hun kansen op; ga na dat de \omterm{def:g11:seq:sequence}{rij} kansen optelt
tot $1$.
\end{method}

\section{Verwachtingswaarde}

\begin{definition}[Verwachtingswaarde]\label{def:g11:prob:expectation}
De \emph{verwachtingswaarde}\index{verwachtingswaarde} (of het \omterm{def:g11:stat:mean}{gemiddelde})
van $X$ is het \omterm{def:g11:stat:mean}{gemiddelde} van haar waarden, gewogen met hun kansen:
\[
\E(X) = \sum_{i=1}^{k} \P(X = x_i)\; x_i .
\]
\end{definition}

\begin{remark}
$\E(X)$ is wat het \omterm{def:g11:stat:mean}{gemiddelde} van een enorm aantal onafhankelijke
herhalingen van het experiment zou zijn --- de precieze versie van die
uitspraak is de wet van de grote aantallen, bewezen in \cref{ch:g12:sums}.
Merk de formele gelijkenis op met het \omterm{def:g11:stat:mean}{gemiddelde} van een gegevensreeks met
\omterm{def:g10:stats:series}{frequenties} (\cref{def:g11:stat:mean}): de kansen spelen de rol van de
relatieve \omterm{def:g10:stats:series}{frequenties} $\frac{n_i}{n}$.
\end{remark}

\begin{example}\label{ex:g11:prob:gameexp}
Voor het spel van \cref{ex:g11:prob:game}:
\[
\E(G) = \frac46 \times (-2) + \frac16 \times 3 + \frac16 \times 4
= \frac{-8 + 3 + 4}{6} = -\frac16 \approx -0.17 .
\]
Gemiddeld verliest de speler ongeveer $17$ cent per spel: het spel is
ongunstig. Een spel heet \emph{eerlijk} wanneer de verwachte winst $0$ is.
\end{example}

\begin{proposition}[Lineariteit]\label{prop:g11:prob:linearity}
Voor alle reële $a, b$ geldt
\[
\E(aX + b) = a\,\E(X) + b .
\]
\end{proposition}

\begin{proof}
De variabele $aX + b$ neemt de waarde $a x_i + b$ aan met \omterm{def:g10:proba:distribution}{kans}
$\P(X = x_i)$, dus
\[
\E(aX + b) = \sum_i \P(X = x_i)(a x_i + b)
= a \sum_i \P(X = x_i)\,x_i + b \underbrace{\sum_i \P(X =
x_i)}_{=\,1} = a\,\E(X) + b .
\]
\end{proof}

\section{Variantie en standaardafwijking}

\begin{definition}[Variantie]\label{def:g11:prob:variance}
De \emph{variantie}\index{variantie} van $X$ meet de spreiding van haar
waarden rond de \omterm{def:g11:prob:expectation}{verwachtingswaarde} $m = \E(X)$:
\[
\V(X) = \E\bigl((X - m)^2\bigr)
= \sum_{i=1}^{k} \P(X = x_i)\,(x_i - m)^2,
\qquad
\sigma(X) = \sqrt{\V(X)} .
\]
\end{definition}

\begin{proposition}[Rekenformule]\label{prop:g11:prob:konig}
\[
\V(X) = \E(X^2) - \E(X)^2 .
\]
\end{proposition}

\begin{proof}
Werk $(x_i - m)^2 = x_i^2 - 2m\,x_i + m^2$ binnen de som uit:
\[
\V(X) = \sum_i \P(X = x_i)\,x_i^2
- 2m \sum_i \P(X = x_i)\,x_i + m^2\sum_i \P(X = x_i)
= \E(X^2) - 2m^2 + m^2 ,
\]
en dat is $\E(X^2) - m^2$. Dit is de versie voor \omterm{def:g11:prob:rv}{toevalsvariabelen} van de
identiteit die in \cref{prop:g11:stat:shortcut} voor gegevensreeksen bewezen
werd.
\end{proof}

\begin{proposition}[Affiene transformatie]\label{prop:g11:prob:affine}
$\V(aX + b) = a^2\,\V(X)$ en $\sigma(aX + b) = \abs a\,\sigma(X)$.
\end{proposition}

\begin{proof}
Door de lineariteit wijkt $aX + b$ van zijn \omterm{def:g11:prob:expectation}{verwachtingswaarde} $a m + b$ af
met $(aX + b) - (am + b) = a(X - m)$: elke afwijking wordt met $a$
geschaald, dus elke gekwadrateerde afwijking met $a^2$, en dus ook hun
gewogen \omterm{def:g11:stat:mean}{gemiddelde}. De verschuiving $b$ is verdwenen --- een variabele
verschuiven verandert haar spreiding niet.
\end{proof}

\begin{example}\label{ex:g11:prob:variance}
Voor het spel van \cref{ex:g11:prob:game} is
$\E(G^2) = \frac46 \times 4 + \frac16 \times 9 + \frac16 \times 16
= \frac{16 + 9 + 16}{6} = \frac{41}{6}$, dus
\[
\V(G) = \frac{41}{6} - \left(-\frac16\right)^2
= \frac{41}{6} - \frac{1}{36} = \frac{245}{36} \approx 6.8,
\qquad
\sigma(G) \approx 2.6 .
\]
Het \omterm{def:g11:stat:mean}{gemiddelde} verlies van $0.17$ per spel gaat gepaard met schommelingen
van typische grootte $2.6$ euro --- een \omterm{def:g11:prob:expectation}{verwachtingswaarde} zonder
\omterm{def:g11:stat:variance}{standaardafwijking} vertelt maar de helft van het verhaal.
\end{example}

\begin{omfigure}
\begin{tikzpicture}
\begin{axis}[
  width=8.6cm, height=5.2cm,
  ybar, bar width=14pt,
  xlabel={$g$}, ylabel={$\P(G = g)$},
  xmin=-3.5, xmax=5.5, ymin=0, ymax=0.8,
  xtick={-2,3,4}, ytick={1/6, 4/6},
  yticklabels={$\tfrac16$,$\tfrac46$},
  tick label style={font=\small},
  label style={font=\small}]
  \addplot[fill=omDef!40, draw=omDef] coordinates {
    (-2,0.6667) (3,0.1667) (4,0.1667)};
  \draw[omThm, dashed, thick] (axis cs:-0.1667,0) --
    (axis cs:-0.1667,0.75);
  \node[omThm, font=\small, anchor=south west] at (axis cs:-0.05,0.62)
    {$\E(G)$};
\end{axis}
\end{tikzpicture}

{\small De \omterm{def:g11:prob:rv}{kansverdeling} van de winst $G$ uit \cref{ex:g11:prob:game}, en
haar \omterm{def:g11:prob:expectation}{verwachtingswaarde} $\E(G) = -\frac16$ (stippellijn): het evenwichtspunt
van de kansmassa's.}
\end{omfigure}

\section{Oefeningen}

\begin{exercise}[$\star$]\label{exo:g11:prob:1}
Er wordt een kaart getrokken uit een standaardspel van $52$ kaarten. Bereken
de \omterm{def:g10:proba:distribution}{kans} om te trekken: een harten; een heer; een harten of een heer; noch
een harten noch een heer.
\end{exercise}

\begin{exercise}[$\star$]\label{exo:g11:prob:2}
Er wordt met twee eerlijke dobbelstenen gegooid en hun som $S$ wordt
genoteerd. Bereken met het model met $36$ uitkomsten uit
\cref{ex:g11:prob:twodice} de kansen $\P(S = 6)$, $\P(S \geq 10)$ en
$\P(S \text{ is even})$.
\end{exercise}

\begin{exercise}[$\star$]\label{exo:g11:prob:3}
Een \omterm{def:g11:prob:rv}{toevalsvariabele} heeft de \omterm{def:g11:prob:rv}{kansverdeling}
\begin{center}
\begin{tabular}{c|cccc}
$x$ & $-1$ & $0$ & $2$ & $5$\\
\hline
$\P(X=x)$ & $0.3$ & $0.2$ & $0.4$ & $p$
\end{tabular}
\end{center}
Zoek $p$, en bereken daarna $\E(X)$.
\end{exercise}

\begin{exercise}[$\star$]\label{exo:g11:prob:4}
$X$ voldoet aan $\E(X) = 3$ en $\V(X) = 4$. Geef de \omterm{def:g11:prob:expectation}{verwachtingswaarde}, de
\omterm{def:g11:prob:variance}{variantie} en de \omterm{def:g11:stat:variance}{standaardafwijking} van $Y = 2X - 5$.
\end{exercise}

\begin{exercise}[$\star\star$]\label{exo:g11:prob:5}
Een rad van fortuin heeft $10$ gelijke sectoren: $5$ tonen $0$ euro, $3$
tonen $2$ euro en $2$ tonen $5$ euro. Meespelen kost $2$ euro. Geef de
\omterm{def:g11:prob:rv}{kansverdeling} van de winst $G$ (uitbetaling min inzet), bereken $\E(G)$, en
beslis of het spel gunstig is.
\end{exercise}

\begin{exercise}[$\star\star$]\label{exo:g11:prob:6}
Een urne bevat $3$ rode en $2$ blauwe ballen; er worden twee ballen
getrokken \emph{zonder} teruglegging. Zij $X$ het aantal getrokken rode
ballen. Bouw de \omterm{def:g11:prob:rv}{kansverdeling} van $X$ (noem de $20$ geordende trekkingen op,
of gebruik een \omterm{met:g10:proba:tree}{boomdiagram}), en bereken $\E(X)$.
\end{exercise}

\begin{exercise}[$\star\star$]\label{exo:g11:prob:7}
Bereken voor het rad van \cref{exo:g11:prob:5} de waarden $\V(G)$ en
$\sigma(G)$.
\end{exercise}

\begin{exercise}[$\star\star$]\label{exo:g11:prob:8}
Een verzekeringsmaatschappij verkoopt een polis voor $80$ euro. Met \omterm{def:g10:proba:distribution}{kans}
$0.05$ dient de klant een schadeclaim van $1000$ euro in, met \omterm{def:g10:proba:distribution}{kans} $0.01$
een van $5000$ euro, en anders niets. Zij $B$ de winst van de maatschappij
op één polis. Geef de \omterm{def:g11:prob:rv}{kansverdeling} van $B$, bereken $\E(B)$, en duid het
resultaat.
\end{exercise}

\begin{exercise}[$\star\star$]\label{exo:g11:prob:9}
Een meerkeuzevraag heeft $4$ antwoorden, waarvan er één juist is. Een juist
antwoord levert $+1$ op; om gokken te ontmoedigen levert een fout antwoord
$x < 0$ op. Een leerling antwoordt willekeurig. Voor welke waarde van $x$ is
de verwachte score van de leerling nul?
\end{exercise}

\begin{exercise}[$\star\star$]\label{exo:g11:prob:10}
Er wordt met twee eerlijke munten gegooid. Ann stelt voor: ``vallen de twee
munten gelijk, dan betaal ik je $1$ euro; verschillen ze, dan betaal jij mij
$1$ euro.'' Bereken de verwachte winst van elke speler. Is het spel
eerlijk? Dezelfde vraag met drie munten, waarbij Ann $2$ euro betaalt als
alle drie gelijk vallen.
\end{exercise}

\begin{exercise}[$\star\star\star$]\label{exo:g11:prob:11}
Een loterij verkoopt $1000$ loten van $5$ euro. Eén lot wint $2000$ euro,
vijf loten winnen $200$ euro, en twintig loten winnen $25$ euro.
\begin{enumerate}
  \item Geef de \omterm{def:g11:prob:rv}{kansverdeling} van de totale opbrengst min de prijzen voor de
        organisator, en van de winst $G$ van één speler.
  \item Bereken $\E(G)$ en de totale verwachte winst van de organisator. Ga
        na dat de twee antwoorden met elkaar kloppen.
\end{enumerate}
\end{exercise}

\section{Opgave: het onderbroken spel}

\begin{problem}[{Weekendopgave --- Pascal, Fermat en de eerlijke verdeling
van een afgebroken weddenschap: de geboorte van de verwachtingswaarde, de
magie van de lineariteit, en wat de verwachtingswaarde niet ziet}]
\label{pb:g11:prob:1}
In 1654 stelde de chevalier de Méré --- dezelfde gokker wiens
dobbelweddenschappen de kansopgave in het onderbouwvolume openden --- aan
Blaise Pascal een diepere vraag: twee even sterke spelers moeten hun reeks
staken bij de stand $5$ tegen $3$, terwijl zes overwinningen de pot van $64$
pistolen opleveren. \emph{Hoe verdeel je die pot eerlijk?} De brieven tussen
Pascal en Fermat waarin ze die vraag beantwoordden, stichtten het begrip dat
in het hart van dit hoofdstuk staat: de \omterm{def:g11:prob:expectation}{verwachtingswaarde}
(\cref{def:g11:prob:expectation}). Je zult hun antwoord op beide manieren
opnieuw afleiden, en daarna kennismaken met de salontruc van de lineariteit
en met de ene blinde vlek van de \omterm{def:g11:prob:expectation}{verwachtingswaarde}.

\medskip
\noindent\textbf{Deel I --- Vlot met de \omterm{def:g11:prob:expectation}{verwachtingswaarde}.}
\begin{enumerate}
  \item Gooi met een eerlijke dobbelsteen; je winst is (ogenaantal $- 3$)
        euro. Bouw de kansverdelingstabel van de winst $G$
        (\cref{met:g11:prob:table}) en bereken $\E(G)$.
  \item Bereken $V(G)$ met de rekenformule (\cref{prop:g11:prob:konig}) en
        $\sigma(G)$ (twee decimalen).
  \item Geef zonder nieuwe tabel de waarden $\E(2G - 1)$ en $V(2G - 1)$
        (\cref{prop:g11:prob:affine}).
  \item Een kraam vraagt $m$ euro om met twee dobbelstenen te gooien en keert
        je de som in euro's uit. Wat is de \emph{eerlijke} prijs $m$? De
        kraam vraagt $8$: wat is je verwachte verlies per spel?
  \item Een laptop van $800$ euro heeft elk jaar $5\,\%$ \omterm{def:g10:proba:distribution}{kans} om vernield te
        worden. Bereken de eerlijke verzekeringspremie, en de verwachte winst
        van de verzekeraar per polis wanneer hij $60$ euro vraagt.
\end{enumerate}

\medskip
\noindent\textbf{Deel II --- De pot verdelen.}
Stand $5$--$3$, wie het eerst $6$ haalt wint de $64$ pistolen, en elke ronde
is een eerlijke muntworp.
\begin{enumerate}[resume]
  \item Twee verleidelijke foute antwoorden: verdelen in de verhouding
        $5 : 3$ (evenredig met de gewonnen rondes), of alles aan de leider
        geven. Zeg in telkens één zin wat er oneerlijk aan is.
  \item De methode van Fermat: speler A heeft nog $1$ overwinning nodig, B
        nog $3$; stel je voor dat er hoe dan ook precies $3$ rondes worden
        gespeeld. Noem de $8$ even waarschijnlijke uitkomsten, duid in elke
        aan wie de reeks wint, en leid de eerlijke verdeling van de $64$
        pistolen af.
  \item Het subtiele punt dat Fermat moest verdedigen: het echte spel zou al
        na één ronde kunnen stoppen --- waarom is het dan toch geoorloofd om
        over alle $3$ denkbeeldige rondes te redeneren?
  \item De methode van Pascal, van achteren naar voren: bereken het eerlijke
        aandeel van A bij de (hypothetische) standen $5$--$5$, daarna
        $5$--$4$, daarna $5$--$3$, en middel telkens de twee even
        waarschijnlijke toekomsten. Vergelijk met vraag 7.
  \item Dezelfde vraag bij de stand $4$--$3$: hoeveel denkbeeldige rondes,
        welke uitkomsten geven B de reeks, en wat is de eerlijke verdeling?
  \item Formuleer het principe dat de twee methoden delen --- de eerlijke
        waarde van een onzekere toekomst is haar \emph{\omterm{def:g11:prob:expectation}{verwachtingswaarde}}
        --- en noem twee moderne bedrijfstakken die er volledig op draaien
        (vraag 5 toont er een).
  \item Een controle voor elke stand: leg met de lineariteit van de
        \omterm{def:g11:prob:expectation}{verwachtingswaarde} (\cref{prop:g11:prob:linearity}) uit waarom de
        eerlijke aandelen van de twee spelers altijd precies de pot van $64$
        pistolen moeten opleveren.
\end{enumerate}

\medskip
\noindent\textbf{Deel III --- De salontruc van de lineariteit.}
\begin{enumerate}[resume]
  \item Haal $\E(\text{som van twee dobbelstenen}) = 7$ in één regel terug
        met de lineariteit (\cref{prop:g11:prob:linearity}), zonder tabel van
        $36$ vakjes. Voor de \emph{\omterm{def:g11:prob:variance}{variantie}} van de som heeft de regel
        $V(X + Y) = V(X) + V(Y)$ nodig dat de dobbelstenen onafhankelijk zijn
        (aangenomen): bereken ze.
  \item Het hoedenfeest: $n$ gasten gooien hun hoed op een stapel en nemen er
        elk willekeurig één terug. Zij $X$ het aantal gasten dat zijn eigen
        hoed terugkrijgt. Wat is voor één vaste gast de \omterm{def:g10:proba:distribution}{kans} om zijn eigen
        hoed terug te krijgen? Tel die $n$ kansen op (lineariteit!) en
        bereken $\E(X)$. Het antwoord zou je moeten verrassen: het hangt niet
        van $n$ af.
  \item Ga $\E(X) = 1$ met brute kracht na voor $n = 2$ en $n = 3$ (noem de
        $2$ en de $6$ even waarschijnlijke hoedentoewijzingen op en middel de
        aantallen).
  \item De terugvondsten van de gasten zijn allesbehalve onafhankelijk (als
        $n - 1$ gasten hun eigen hoed terugkregen, moet de laatste dat ook).
        Waar had de berekening van vraag 14 dan precies \emph{geen}
        onafhankelijkheid nodig? Formuleer de superkracht van de lineariteit.
\end{enumerate}

\medskip
\noindent\textbf{Deel IV --- Wat de \omterm{def:g11:prob:expectation}{verwachtingswaarde} niet ziet.}
\begin{enumerate}[resume]
  \item Twee kraskaarten van $5$ euro: kaart A wint $10$ euro met \omterm{def:g10:proba:distribution}{kans}
        $\frac12$; kaart B wint $1\,000$ euro met \omterm{def:g10:proba:distribution}{kans} $0.005$. Ga na dat
        beide winsten dezelfde \omterm{def:g11:prob:expectation}{verwachtingswaarde} hebben, en bereken daarna
        beide \omterm{def:g11:prob:variance}{varianties}. Welke kaart kopen echte mensen, en wat kopen ze
        eigenlijk?
  \item Het Sint-Petersburgspel: gooi tot de eerste keer munt; valt de eerste
        munt bij worp $n$, dan ontvang je $2^n$ euro. Bereken de bijdrage van
        elke $n$ aan de \omterm{def:g11:prob:expectation}{verwachtingswaarde}, en de \omterm{def:g11:prob:expectation}{verwachtingswaarde} zelf.
        Zou je $1\,000$ euro betalen om één keer te spelen? Wat leert die
        botsing over de \omterm{def:g11:prob:expectation}{verwachtingswaarde}?
  \item Verzekering vanuit de klant bekeken: $20$ euro betalen, of een risico
        van $1\,\%$ lopen om $1\,000$ te verliezen? Vergelijk de twee
        \omterm{def:g11:prob:expectation}{verwachtingswaarden}, en leg uit waarom de verzekering kopen toch
        verstandig kan zijn (waarover kan één gezin niet ``middelen''?).
  \item Slotstuk --- de levensloop van de \omterm{def:g11:prob:expectation}{verwachtingswaarde}, één zin per
        hoofdstuk: geboren in 1654 om een pot te verdelen; haar superkracht,
        de lineariteit, die \omterm{def:g11:prob:expectation}{verwachtingswaarden} optelt zonder naar
        onafhankelijkheid te vragen (de hoeden); haar blinde vlek, het
        risico, dat de \omterm{def:g11:prob:variance}{variantie} meet (de kraskaarten, Sint-Petersburg, de
        verzekering); en haar kroning volgend jaar, wanneer de wet van de
        grote aantallen precies uitlegt waarom de bank altijd wint.
\end{enumerate}
\end{problem}
